\chapter{Completion of the proof of Theorem~{\ref{MAIN}}}

\section{Proof of the strong canonical neighborhood assumption}
\label{sectcannbhd}

\begin{proposition}[The strong canonical neighborhood assumption extends to the next scale $r_{i+1}$]
\label{prop:canonical-neighborhood-extends}
\uses{prop:kappa-limit-noncollapsing, thm:main-existence-ricci-flow-surgery}

\label{extend}
Suppose positive nonincreasing sequences
\[
  \delta_0\geq\cdots\geq\delta_i,
  \qquad
  \epsilon=r_0\geq r_1\geq\cdots\geq r_i,
  \qquad
  \kappa_0\geq\cdots\geq\kappa_i
\]
are fixed. For every $r_{i+1}\leq r_i$, let $\delta(r_{i+1})$ be the constant
from Proposition~\ref{KAPPALIMIT}. Then one can choose
$r_{i+1}\leq r_i$ and $\delta_{i+1}\leq\delta(r_{i+1})$ such that a flow on
$[0,T)$, $T\in(T_i,T_{i+1}]$, satisfying the standing assumptions and the
inductive curvature, noncollapsing, and canonical-neighborhood hypotheses before
$T_i$, with $\overline\delta(t)\leq\delta_{i+1}$ on
$[T_{i-1},T]$, satisfies the strong $(C,\epsilon)$-canonical-neighborhood
assumption with parameter $r_{i+1}$ on $[0,T)$.
\end{proposition}

\begin{proof}
  \uses{prop:kappa-limit-noncollapsing, kaplimit, limitcannbhd, sliceconv, canonvary, narrows, controlnbhd, thm:local-surgery, prop:long-L-length-near-surgery-cap, lem:surgery-cap-ball-bounds, defncanonnbhd, defnepsround, Ccomponent, neckcurv, A_0, Restim, stdsoln, bigell, neckglue, stdsolncannbhd, def:overline-delta-0, prop:surgery-cap-evolution-alternative}
\sketch

Assume failure for every choice and take $r_a\to0$, then surgery controls
$\delta_{a,b}\to0$. Lemma~\ref{16.2} produces first failure points at the
scale $r_a$. Rescale each flow at such a point to scalar curvature one. The
rescaled flows are uniformly noncollapsed and curvature-pinched. The failure
sequence either has uniformly long backward flow lines, yielding a smooth
ancient limit and hence a $\kappa$-solution, or reaches a surgery cap in bounded
backward time. The latter is excluded by Lemma~\ref{lemstdcan}; the former is
excluded by the canonical-neighborhood description of $\kappa$-solutions. This
contradiction proves the extension.
\end{proof}

\begin{lemma}[A limiting counterexample sequence of failure points converges to a limit failure point]
\label{lem:limit-canonical-neighborhood-fails-at-surgery-scale}
\label{16.2}
For each $a$, and all sufficiently large $b$, there is a time
$t_{(a,b)}\in[T_i,T_{(a,b)})$ such that the flow before $t_{(a,b)}$ satisfies
the strong $(C,\epsilon)$-canonical-neighborhood assumption with parameter
$r_a$, while at some point $x_{(a,b)}$ in the slice at $t_{(a,b)}$ the assumption
fails.
\end{lemma}

\begin{proof}
\sketch
\uses{thm:local-surgery,prop:long-L-length-near-surgery-cap,lem:surgery-cap-ball-bounds,defncanonnbhd,defnepsround,Ccomponent,canonvary,neckcurv,A_0,Restim,stdsoln,bigell,neckglue,stdsolncannbhd}
Choose a sequence of counterexample points with times decreasing to the first
failure time. Compactness of the preceding time-slice gives a limiting point
$x_{(a,b)}$ with $R\geq r_a^{-2}$. Stability under smooth time convergence
preserves round components, $C$-components, and cap cores. If the limiting
neighborhood is a strong neck, either its backward embedding continues through
the relevant interval, in which case nearby points inherit strong necks, or it
meets a surgery cap. In the latter case rescale by the cap scale; the standard
solution estimates and neck gluing produce a canonical cap around nearby failure
points, again a contradiction. Thus the first failure point exists as stated.
\end{proof}

\begin{remark}[The limiting canonical neighborhood is a cap, not a strong epsilon-neck]
\label{rem:limit-canonical-neighborhood-is-cap}
\uses{lem:limit-canonical-neighborhood-fails-at-surgery-scale}

When the limiting strong-neck embedding meets a surgery cap, the canonical
neighborhood obtained for the approximating points is a $(C,\epsilon)$-cap
coming from the cap evolution, not a strong $\epsilon$-neck.
\end{remark}

\begin{claim}[The rescaled surgery time lies before $\theta_1$ in the limit]
\label{claim:tilde-t-prime-bounded-by-theta1}
\uses{def:overline-delta-0}

After rescaling a counterexample by the previous surgery scale and moving that
surgery time to $0$, there is $\theta_1<1$, independent of the sequence, such
that the rescaled failure time $\widetilde t'\leq\theta_1$.
\end{claim}

\begin{proof}
\sketch
\uses{prop:surgery-cap-evolution-alternative,def:overline-delta-0}
Embed the largest standard-solution ball backward from the cap. If the embedding
ends before $\theta_1$, the failure time is earlier. If it reaches $\theta_1$,
the standard solution has scalar curvature at least $8D$, while the neck flow
line has curvature at most $6D$; the cap-evolution comparison contradicts this.
\end{proof}

\begin{claim}[The backward-flowed base point lies in the rescaled standard-flow ball]
\label{claim:w-ab-in-ball-image}
For sufficiently small surgery control, the point obtained by flowing a
rescaled counterexample backward to the surgery time lies in the image of the
standard-flow ball $B(p_0,A)$.
\end{claim}

\begin{proof}
\sketch
The strong-neck coordinate gives distance at most
$1.1\widetilde Q^{-1/2}\epsilon^{-1}$ from the backward point to the cap point,
and the cap-ball bound gives distance at most $A_0+5$ from that point to the
cap tip. The curvature comparison
\begin{equation}
\label{DQ}
  (5D)^{-1}\leq\widetilde Q\leq5D
\end{equation}
and the chosen radius $A$ put the backward point inside the standard-flow ball.
\end{proof}

\begin{claim}[Limit points of scalar curvature exceeding one have a doubled canonical neighborhood]
\label{claim:limit-points-have-canonical-neighborhood}
For all sufficiently large indices, every point in the rescaled zero-time slice
with $R>1$ has a $(2C,2\epsilon)$-canonical neighborhood.
\end{claim}

\begin{proof}
\sketch
\uses{thm:local-surgery,omegacanon,epslimit}
Exposed points lie in standard-solution cap cores when the surgery control tends
to zero. At nonexposed points, approach from negative times, where the original
canonical-neighborhood assumption applies. Smooth stability preserves round
components, $C$-components, and cap cores; the neck-limit result preserves a
strong neck with parameters $(2C,2\epsilon)$.
\end{proof}

\begin{lemma}[A bounded-curvature flow line into a surgery cap forces a canonical neighborhood]
\label{lem:std-cap-canonical-neighborhood}
\label{lemstdcan}
Suppose that, for fixed $A',D',t'$, each rescaled flow has a point in
$B(\widetilde x_a,0,A')$ whose backward flow line reaches a surgery cap within
time $t_a\leq t'$ and has scalar curvature at most $D'$ along the line. Then
for all sufficiently large $a$, $x_a$ has a strong $(C,\epsilon)$-canonical
neighborhood.
\end{lemma}

\begin{proof}
  \uses{A_0, stdsoln, prop:surgery-cap-evolution-alternative, posdistdec, stdsolncannbhd, canonvary, neckglue, def:overline-delta-0, Restim}
\sketch

The cap curvature bound gives a uniform lower bound
\begin{equation}
\label{hequat}
  \bar h_a^2\geq(2D'D)^{-1}
\end{equation}
for the surgery scale. Claim~\ref{claim:elapsed-time-bounded-by-theta1} bounds
$t_a/\bar h_a^2<1$. Standard-solution compactness then places a uniformly large
ball around $x_a$ inside the evolved cap. The standard canonical-neighborhood
theorem gives either a cap or an evolving neck. A cap transfers directly by
smooth convergence; in the neck case, its initial slice lies in the continuing
region and Proposition~\ref{neckglue} glues it to the surgery neck. Both cases
contradict failure of the canonical-neighborhood assumption.
\end{proof}

\begin{claim}[The elapsed time to the surgery cap is bounded by $\theta_1$ times the surgery scale squared]
\label{claim:elapsed-time-bounded-by-theta1}
There is $\theta_1<1$, depending only on $D'$ and $t'$, such that
$t_a<\theta_1\bar h_a^2$ for all sufficiently large $a$.
\end{claim}

\begin{proof}
\sketch
\uses{prop:surgery-cap-evolution-alternative,def:overline-delta-0,Restim}
If $t_a\leq\bar h_a^2/2$, choose any $\theta_1\in(1/2,1)$. Otherwise the
bounded-curvature hypothesis bounds the cap flow line by
$2t'D'\bar h_a^{-2}$. Choose $\theta_1$ so that the standard solution has
curvature at least $6t'D'$ after rescaled time $2\theta_1-1$. The cap comparison
then contradicts the assumed upper bound if $t_a\geq\theta_1\bar h_a^2$.
\end{proof}

\begin{claim}[Uniform curvature bound along backward flow lines from a fixed ball]
\label{claim:curvature-bound-along-backward-flow-lines}
\label{D(A)delta(A)}
For every $A<\infty$ there are $D(A)<\infty$ and $\delta(A)>0$ such that,
for all sufficiently large $a$, $|\Rm|\leq D(A)$ along every backward flow line
from $B(\widetilde x_a,0,A)$ defined for time at most $\delta(A)$.
\end{claim}

\begin{proof}
\sketch
\uses{prop:kappa-limit-noncollapsing,bcbd,controlnbhd}
The rescaled flows are uniformly noncollapsed and points of scalar curvature
above one have doubled canonical neighborhoods. The bounded-curvature-at-
bounded-distance estimate bounds scalar curvature on the fixed ball; the
canonical-neighborhood evolution inequality then extends this bound for a
uniform short backward time. Positive pinching converts the scalar bound to a
full curvature bound.
\end{proof}

\begin{claim}[Dichotomy: uniform backward parabolic neighborhood or a canonical neighborhood]
\label{claim:parabolic-neighborhood-or-canonical-neighborhood}
After passing to a subsequence, either for every $A$ there are $D(A),t(A)>0$
such that $P(\widetilde x_a,0,A,-t(A))$ exists with $|\Rm|\leq D(A)$, or every
$x_a$ has a strong $(C,\epsilon)$-canonical neighborhood.
\end{claim}

\begin{proof}
\sketch
\uses{claim:curvature-bound-along-backward-flow-lines,lem:std-cap-canonical-neighborhood}
If a uniform backward parabolic neighborhood exists, Claim~\ref{D(A)delta(A)}
supplies the curvature bound. Otherwise a bounded-radius ball contains a flow
line reaching a surgery cap in arbitrarily short time; the same claim bounds its
curvature, and Lemma~\ref{lemstdcan} gives a canonical neighborhood.
\end{proof}

\begin{claim}[If the backward time bound cannot be made uniform in the radius, canonical neighborhoods exist]
\label{claim:uniform-t-prime-independent-of-A}
If $t(A)$ in Claim~\ref{claim:parabolic-neighborhood-or-canonical-neighborhood}
cannot be chosen independent of $A$, then a subsequence of the $x_a$ has
strong $(C,\epsilon)$-canonical neighborhoods.
\end{claim}

\begin{proof}
\sketch
\uses{lem:std-cap-canonical-neighborhood,controlnbhd}
The zero-time slices converge on every fixed ball to a complete bounded
nonnegatively curved limit. A uniform short-time curvature estimate for this
limit either extends the backward parabolic neighborhoods uniformly in $A$ or
produces a bounded-curvature flow line into a surgery cap. The latter invokes
Lemma~\ref{lemstdcan}, so the assumed nonuniformity is impossible.
\end{proof}

The preceding dichotomy and the last claim give a backward time $T'>0$ uniform
over all radii. If $T'=\infty$, the rescaled sequence converges to a
$\kappa$-solution, whose canonical-neighborhood classification contradicts the
choice of $x_a$. If $T'<\infty$, maximality produces a flow line into a surgery
cap with uniformly bounded curvature, again contradicting
Lemma~\ref{lemstdcan}. This proves Proposition~\ref{extend}.

\section{Surgery times do not accumulate}
\label{sectcomplete}

Choose $r_{i+1},\delta_{i+1}$ by Proposition~\ref{extend} and set
$\kappa_{i+1}=\kappa(r_{i+1})$ from Proposition~\ref{KAPPALIMIT}. If a flow
becomes singular at $T<T_{i+1}$, perform the surgery operation at $T$ and
restart the flow from the compact post-surgery slice. Lemma~\ref{14.26} and
Proposition~\ref{extend} preserve all inductive properties. If the controlled
limiting region is empty, the flow becomes empty and the induction terminates.
Otherwise repeat at the next singular time.

\begin{lemma}[The number of surgery times up to $T_{i+1}$ is bounded]
\label{lem:bounded-number-of-surgeries}
\uses{thm:main-existence-ricci-flow-surgery}

\label{bddsurgery}
For a flow satisfying Theorem~\ref{MAIN} on $[0,T)$ with $T\leq T_{i+1}$,
the number of surgery times is bounded by a constant depending only on
$\Vol(M_0,g(0))$, $T_{i+1}$, $r_{i+1}$, and
$\overline\delta(T_{i+1})$.
\end{lemma}

\begin{proof}
\sketch
\uses{A_0}
Away from surgery, positive pinching gives
\[
  \frac{d}{dt}\Vol(M_t)=-\int_{M_t}R\,d\vol\leq6\Vol(M_t),
\]
so volume grows at most exponentially. Each neck surgery removes a definite
volume: the removed horn contains a positive half-neck of scale
$h(t)\geq h(T_{i+1})>0$, while the inserted standard ball has volume bounded by
$(1+\epsilon)Kh(t)^3$. Thus every surgery decreases volume by a fixed amount
depending only on the listed parameters. The number of $2$-sphere surgeries is
bounded; the number of components is bounded by the initial number plus the
number of surgeries, and normalized initial volume bounds the initial number of
components. Hence the total number of surgeries is bounded.
\end{proof}

Therefore surgery times are finite on every finite interval, completing the
inductive proof of Theorem~\ref{MAIN}.
